\makeabstract
{
Estimating Enthalpies of Solution Using Infrared Thermography and Coffee Cup Calorimetry
}
{
Aidan Cahill (1), James Shriver (2), Dorothy Huang (3)
}
{
\textsuperscript{1}Department of Chemistry, Amherst College, Amherst, MA, USA
}
{
To create a general chemistry lab featuring infrared thermography, we used the FLIR ONE camera to estimate the enthalpy of solution of the following ionic solids: sodium chloride, sodium bromide, lithium chloride, lithium bromide, potassium chloride, potassium bromide, and ammonium chloride. 
}
{
By measuring the difference in temperature before and after the solute was dissolved, we determined if the reaction was endothermic or exothermic. Furthermore, we used the formula -∆Hsol = qH2O = mc∆T (with Ccal accounted for) to estimate the enthalpy of solution, which we subsequently compared to the literature values.
}
{
We found potassium chloride, potassium bromide, and ammonium chloride to be significantly endothermic; sodium chloride and sodium bromide to be near zero; and lithium chloride and lithium bromide to be significantly exothermic. The observed trend likely occurred because ions of smaller atomic radii have stronger interactions with water, resulting in greater ion-dipole interactions and more exothermic enthalpies of solution. All such values roughly correspond to the literature values, though the magnitude of their estimated enthalpy of solution was reduced by about three times. For more accurate measurements in the future, the specific heat capacity of the solutions can be found using a differential scanning calorimeter (DSC).
}
{
Overall, merging coffee cup calorimetry with infrared thermography offers an effective way to estimate enthalpies of solution and visualize heat transfer. We conclude that this experiment would successfully educate general chemistry students on thermodynamics-related concepts.  

}

\newpage
